Let $B>1$, $\ell\ge B$ and $d=d_n(\ell)\ge f_n(B)$, where $f_n(B)=\max\{1,\,1+\lfloor\log C_n/\log B\rfloor\}$ and $C_n>0$ (Prop.~D: $D_n>0$).

\emph{The integer floor.} Write $x=\log C_n/\log B\in\mathbb R$ (the sign of $x$ is irrelevant; for $C_n<1$ it is negative). By definition of the floor, $\lfloor x\rfloor\le x<\lfloor x\rfloor+1$, so $f_n(B)\ge\lfloor x\rfloor+1>x$; this is the only property of $f_n(B)$ used, and it is why the ``$+1$'' is in the definition: $\lfloor x\rfloor+1$ is the least integer strictly larger than $x$. The outer $\max\{1,\cdot\}$ only matters when $C_n<1$ (then $x<0$ and $\lfloor x\rfloor+1\le0$); it keeps $f_n(B)\ge1=\min_\ell d_n(\ell)$, so that the hypothesis $d\ge f_n(B)$ is never vacuous, and it does not affect the inequality $f_n(B)>x$. (Without the $\max$ the statement would still be true, because $d\ge1$ for every odd $\ell$; the $\max$ removes the case distinction.)

\emph{From the floor to the constant.} Since $B>1$, $\log B>0$ and $x\mapsto B^x=\exp(x\log B)$ is strictly increasing on $\mathbb R$. Hence $B^{f_n(B)}>B^{x}=\exp(\log C_n)=C_n$. If $d\ge f_n(B)$ then $B^d\ge B^{f_n(B)}$ (monotonicity again, $B>1$), so $B^d>C_n$. If instead $d<f_n(B)$ is excluded by hypothesis, nothing else is needed: we only use $d\ge f_n(B)$ and $d\ge1$.

\emph{From $B$ to $\ell$.} $\ell\ge B>1$ and $d\ge1$ give $\ell^d\ge B^d$ (monotonicity of $y\mapsto y^d$ on $y>0$ for a positive integer exponent). Therefore $\ell^{d_n(\ell)}=\ell^d\ge B^d>C_n$, which is the hypothesis $\ell^{d_n(\ell)}>C_n$ of Theorem~\ref{thm:A}, and Theorem~\ref{thm:A} gives $\ell\nmid k_n$.

\emph{Example (the numbers of Cor.~7).} At $n=7$ with $B=1314283897427173$: $\log C_7/\log B=\log_{10}C_7/\log_{10}B=30.2374\ldots/15.1187\ldots=1.99998\ldots$ (the ratio does not depend on the base; in natural logarithms $69.6241\ldots/34.8121\ldots$), so $f_7(B)=1+1=2$; every prime $\ell\ge B$ with $d_7(\ell)\ge2$, i.e.\ $\ell\not\equiv1\pmod{128}$, satisfies $\ell^2\ge B^2>C_7$ (the integer form of this comparison is Lemma~\ref{lem:C7int}). With $B=1727342163036353095979941756930$ the quotient is $0.99999\ldots$, $f_7(B)=1$, and every prime $\ell\ge B$ is covered whatever its order. Both thresholds are the frozen integers of Cor.~\ref{cor:7}.

\emph{Lean.} \path{TheoremAInputs.theoremA_threshold} takes the two comparisons $C<B^d$ and $B^d\le\ell^d$ as hypotheses and feeds them to the core \path{theoremA_of_inputs}; the floor computation above is the natural-language justification of those two hypotheses and is not itself formalised (it is ordinary real-number monotonicity).
